% =====================================================================
% abstract.tex — DRAFTED at rb-047 (RB-047), 2026-05-30.
%
% Filled after the body of the manuscript landed:
%   - sec:model (rb-009, RB-004), three-lever reframe.
%   - sec:results-c1..c4 (rb-006, rb-007, rb-011, rb-013).
%   - sec:extensions-A3, A2, A8 (rb-017, rb-022, rb-028).
%   - sec:appendix-c5 (rb-018), sec:appendix-deriv-c2 (rb-024),
%     sec:appendix-deriv-a3 (rb-046).
%
% Strength check (mission §3, CLAIM_LEDGER.md): every quantitative
% claim below is at the licensed band / distributional / conditional
% strength.  Original-paper categorical phrasings (CF floor [0.60,
% 0.96]; "negligible regardless of other parameters"; uniform
% no-inversion; A1 upper bound on VDA) appear here only as the thing
% being retracted, not stated.
% =====================================================================

\begin{abstract}
\noindent
In a cued change-detection task whose cue carries both value
$\val \ge 1$ and validity $\valid \in [1/\Nloc, 1]$, an ideal observer
can adapt to value by shifting decision criteria, by re-allocating
spatial attention, or, when noise is correlated across locations, by
\emph{decorrelating} the encoded population code.  We rebuild the
normative model of \emph{value-directed attention} (VDA) around this
three-lever decomposition (Section~\ref{sec:model},
Definition~\ref{def:three-levers}), with cross-location correlation
$\corr$ promoted to a first-class model parameter via an exact
one-factor Gauss--Hermite reduction that recovers the inherited
independent-noise model in the limit $\corr \to 0$ to floating-point
identity.  The rebuilt manuscript reports every headline result at
its defensible distributional, graded, or conditional strength
rather than as a categorical floor.

Across a $4{,}410$-cell sweep over $(\valid, \val, \Rsens)$, the
criterion fraction $\CF$ is concentrated in $[0.30, 1.00]$ with
median $0.76$; the published $[0.60, 0.96]$ interval is retracted on
both ends, and adding $\corr = 0.2$ drops the strict minimum below
$0.5$ (Section~\ref{sec:results-c1}).  $\VDA$ is non-monotonic in
$\Rsens$ as in the inherited paper; we strengthen this result with
the closed-form escape threshold
$\rdagger(\val) = K_u(\val) / [(\Nloc - 1) K_c(\val)]$ that pins the
lower edge of the VDA-active band (Section~\ref{sec:results-c2}),
extend it to $\corr > 0$ via the natural $\corr$-aware gradient
quadratures (Appendix~\ref{sec:appendix-deriv-c2}), and prove that
$\rdagger(\val)$ is invariant under the conservation order of the
benefit/cost rule (Appendix~\ref{sec:appendix-deriv-a3}).  The
``narrow regime'' for VDA becomes a quantitative iso-$\VDA$ contour
band (Section~\ref{sec:results-c3}): at $\valid \ge 0.95$ peak
$\VDA$ sits at the grid floor under any $(\Rsens, \corr)$; at
$\valid \ge 0.80$ this survives only at $\corr = 0$; at
$\valid \ge 0.60$ it fails.  The ``no inversion'' result is restated
as a conditional theorem with the explicit condition
$\valid \ge 1 / [(\Nloc - 1) \val + 1]$, and the rebuild adds the
\emph{anti-cue inversion} as a new falsifiable prediction
(Section~\ref{sec:results-c4}, Equation~\eqref{eq:value-weight}):
under counter-predictive cues ($\valid < 1/\Nloc$) at $\Nloc = 4$,
the globally optimal allocation places \emph{less} attention at the
high-value location in $36.1\%$ of probed cells.  Symmetric recovery
at $\Rsens = 1$ is reaffirmed as a universal real-number identity
(Appendix~\ref{sec:appendix-c5}, conservation-form-invariant by
construction).

Independence is not an upper bound on $\VDA$.  What independence
actually upper-bounds is the criterion fraction, and only in
variant~A: $\partial \VDA / \partial \corr$ flips sign near
$\Rsens \approx 0.5$, with the cell-wise crossover sweeping past
$\Rsens \approx 0.79$ across the full sweep
(Section~\ref{sec:model-upper-bound}).  Within-display heterogeneity
of $\Rsens$ is a bounded perturbation that leaves the $C_2$ peak
coordinates and the contested-CF corner essentially invariant
(Section~\ref{sec:extensions-a2}); the additive
$\benefit + \cost = 2$ rule is one point in the power-mean
conservation family, and the headline numbers are reported as a
band over that family (Section~\ref{sec:extensions-a3}); the
homogeneous-uncued result is innocuous at the inherited additive
optimum but can fail in a benefit-dominant symmetric-stress corner
when the conservation rule is multiplicative
(Section~\ref{sec:extensions-a8}).  Every quantitative claim above
is backed by a simulation under \texttt{Rebuild/sims/} with a
recovered limit, a deterministic output hash, and a captioned
figure; every novel proposition is derived under
\texttt{Rebuild/derivations/}; nothing in this abstract is stated
more strongly than its row of \texttt{CLAIM\_LEDGER.md} licenses.
\end{abstract}
